\documentclass[11pt,a4paper]{article}
\usepackage[utf8]{inputenc}
\usepackage[T1]{fontenc}
\usepackage[german,english]{babel}
\usepackage{amsmath,amssymb,amsthm}
\usepackage{geometry}
\usepackage{booktabs}
\usepackage{cite}
\usepackage[colorlinks=true,linkcolor=blue,citecolor=blue,urlcolor=blue]{hyperref}

\geometry{verbose,tmargin=2.5cm,bmargin=2.5cm,lmargin=2.5cm,rmargin=2.5cm}

\newtheorem{definition}{Definition}
\newtheorem{proposition}{Proposition}
\newtheorem{theorem}{Satz}
\newtheorem{remark}{Bemerkung}

\title{\textbf{Radiale Initialtopologie, Kolmogorov-Quotient und meromorphe Struktur der $E_8$-Epstein-Zetafunktion}}
\author{\textbf{Thomas Hoffbauer} \\ \small Computer Institut Bamberg (CIB)}
\date{1. August 2026}

\begin{document}
	
	\maketitle
	
	\begin{abstract}
		\selectlanguage{english}
		We present an expository note on the topological and analytic properties of the $E_8$ lattice norm decomposition. We introduce the radial initial topology $\tau_{\text{rad}} = \{Q^{-1}(V) \mid V \subseteq \mathbb{N}_0\}$ generated by the norm quadratic form $Q(x) = \|x\|^2/2$ on the even unimodular $E_8$ lattice. We show that the associated Kolmogorov quotient space $E_8 / \sim$ is canonically homeomorphic to $\mathbb{N}_0$. Utilizing the classical theta identity $\Theta_{E_8} = E_4$, we recall the radial Epstein zeta function $Z_{E_8}(s)$ and its analytic continuation. We confirm that $Z_{E_8}(s)$ possesses exactly one simple pole at $s=8$ with residue $\text{Res}_{s=8} Z_{E_8}(s) = \pi^4 / 3$, whereas the apparent singularity at $s=2$ is removable, yielding the continuation value $Z_{E_8}(2) = -30\zeta(3)/\pi^2$. Finally, we review the asymptotic growth of discrete norm balls and the arithmetic scaling decomposition for primitive lattice vectors.
	\end{abstract}
	
	\selectlanguage{german}
	
	\section{Einleitung und grundlegende Definitionen}
	
	Das gerade unimodulare $E_8$-Gitter $\mathbf{E}_8 \subset \mathbb{R}^8$ nimmt aufgrund seiner herausragenden Symmetrie- und Dichteeigenschaften eine fundamentale Stellung in der diskreten Geometrie und der Theorie der Modulformen ein \cite{serre1973, conway1999}. Die durch das Gitter $\mathbf{E}_8$ bestimmte Kugelpackung ist unter allen Kugelpackungen in $\mathbb{R}^8$ bewiesen optimal \cite{viazovska2017}. In dieser expositorischen Note geben wir eine strukturierte topologische Reformulierung der Schalenzerlegung von $\mathbf{E}_8$ und fassen die analytischen Eigenschaften der zugehörigen radialen Epstein-Zetafunktion präzise zusammen.
	
	Sei $\mathbf{E}_8 \subset \mathbb{R}^8$ definiert durch:
	$$\mathbf{E}_8 = \left\{ x \in \mathbb{Z}^8 \cup \left(\mathbb{Z} + \tfrac{1}{2}\right)^8 \;\middle|\; \sum_{i=1}^8 x_i \in 2\mathbb{Z} \right\}$$
	Wir betrachten die normierte quadratische Form $Q(x) := \frac{\|x\|^2}{2} \in \mathbb{N}_0$ für alle $x \in \mathbf{E}_8$.
	
	\begin{definition}[Schalen-Partition]
		Für $n \in \mathbb{N}_0$ ist die $n$-te $E_8$-Normschale $S_n$ definiert als die Faser von $Q$ über $\{n\}$:
		$$S_n := Q^{-1}(\{n\}) = \left\{ x \in \mathbf{E}_8 \;\middle|\; \frac{\|x\|^2}{2} = n \right\}$$
		Der Gesamtraum zerfällt in die disjunkte Vereinigung $\mathbf{E}_8 = \bigsqcup_{n=0}^{\infty} S_n$.
	\end{definition}
	
	\begin{definition}[Radiale Initialtopologie]
		Die radiale Topologie $\tau_{\text{rad}}$ auf $\mathbf{E}_8$ ist die von der Schalenabbildung $Q: \mathbf{E}_8 \to \mathbb{N}_0$ aus der diskreten Topologie auf $\mathbb{N}_0$ zurückgezogene Initialtopologie:
		$$\tau_{\text{rad}} := \left\{ Q^{-1}(V) \;\middle|\; V \subseteq \mathbb{N}_0 \right\}$$
		Sie ist zugleich die durch die Faserpartition $\mathbf{E}_8 = \bigsqcup_{n \ge 0} S_n$ bestimmte Partitionstopologie.
	\end{definition}
	
	\section{Topologische und ordnungstheoretische Struktur}
	
	\begin{proposition}[Kolmogorov-Quotient]
		Die Funktion $d_{\text{rad}}(x,y) := |Q(x) - Q(y)|$ ist eine Pseudometrik auf $\mathbf{E}_8$ und induziert die radiale Topologie $\tau_{\text{rad}}$, da die Metrik $|m-n|$ auf $\mathbb{N}_0$ die diskrete Topologie erzeugt. Es gilt $d_{\text{rad}}(x,y) = 0 \iff Q(x) = Q(y)$. Insbesondere sind verschiedene Punkte derselben Normschale topologisch ununterscheidbar, sodass $(\mathbf{E}_8, \tau_{\text{rad}})$ nicht $T_0$ ist.
		
		Für die Äquivalenzrelation $x \sim y \iff Q(x) = Q(y)$ ist der Kolmogorov-Quotient $(\mathbf{E}_8 / \sim, \tau_{\text{quot}})$ kanonisch homöomorph zum diskreten Bild $Q(\mathbf{E}_8) \subseteq \mathbb{N}_0$. Wie Satz~\ref{thm:theta} zeigt, gilt $|S_n| = 240\,\sigma_3(n) > 0$ für alle $n \ge 1$, sodass $Q(\mathbf{E}_8) = \mathbb{N}_0$ folgt. Folglich gilt:
		$$\mathbf{E}_8 / \sim \;\;\cong\;\; \mathbb{N}_0$$
		Jede Schale $S_n$ ist unter der Automorphismengruppe $\operatorname{Aut}(\mathbf{E}_8)$ invariant. Da $S_n$ endlich ist, zerfällt sie in endlich viele Automorphismenorbits.
	\end{proposition}
	
	\begin{proposition}[Alexandrov-Radialtopologie]
		Die durch die Präordnung $x \preceq y \iff Q(x) \le Q(y)$ induzierte Alexandrov-Topologie $\tau_{\uparrow}$ besteht aus den Urbildern aufwärts abgeschlossener Teilmengen von $\mathbb{N}_0$. Da abgeschlossene Mengen in der Alexandrov-Topologie genau die nach unten abgeschlossenen Teilmengen bezüglich $\preceq$ sind, ist der diskrete Normball $K_N := \{x \in \mathbf{E}_8 \mid Q(x) \le N\}$ abgeschlossen. Da jede Schale $S_n$ endlich ist und $K_N = \bigsqcup_{n=0}^N S_n$ gilt, ist $K_N$ endlich und folglich kompakt.
	\end{proposition}
	
	\section{Schalenarithmetik und Ballasymptotik}
	
	\begin{theorem}[Theta-Identität und Ballasymptotik]
		\label{thm:theta}
		Nach der klassischen Identität für die $E_8$-Theta-Reihe (Eisensteinreihe $E_4$) gilt \cite{serre1973}:
		$$\Theta_{\mathbf{E}_8}(q) = \sum_{x \in \mathbf{E}_8} q^{Q(x)} = 1 + 240 \sum_{n=1}^{\infty} \sigma_3(n) q^n \implies |S_n| = 240 \, \sigma_3(n) \quad (n \ge 1)$$
		wobei $\sigma_3(n) = \sum_{d|n} d^3$ die Teilerfunktion dritter Potenz bezeichnet \cite{apostol1976}. Da $\sum_{n \le N} \sigma_3(n) \sim \frac{\zeta(4)}{4} N^4$ gilt, folgt für die kumulierte Punktzahl des diskreten Normballs $K_N$ das asymptotische Wachstumsgesetz \cite{apostol1976}:
		$$|K_N| = 1 + 240 \sum_{n=1}^{N} \sigma_3(n) \;\sim\; 240 \cdot \frac{\zeta(4)}{4} N^4 = 60 \cdot \frac{\pi^4}{90} N^4 = \frac{2\pi^4}{3} N^4$$
	\end{theorem}
	
	\section{Analytische Fortsetzung der radialen Epstein-Zetafunktion}
	
	Wir verwenden für die radiale Epstein-Zetafunktion die spektrale Konvention $Z_{\mathbf{E}_8}(s) := \sum_{x \neq 0} \|x\|^{-s}$. Für die durch $Q(x) = \|x\|^2/2$ definierte Standard-Epstein-Funktion $\mathcal{E}_{\mathbf{E}_8}(u) := \sum_{x \neq 0} Q(x)^{-u}$ gilt die Beziehung $\mathcal{E}_{\mathbf{E}_8}(u) = 2^u Z_{\mathbf{E}_8}(2u)$, sodass die Variablen durch $u = s/2$ verbunden sind \cite{terras1985}.
	
	\begin{theorem}[Meromorphe Fortsetzung und Polstellenstruktur]
		Für $\operatorname{Re}(s) > 8$ konvergiert die radiale Epstein-Zetafunktion absolut und erfüllt:
		$$Z_{\mathbf{E}_8}(s) = 240 \cdot 2^{-s/2} \cdot \zeta\left(\frac{s}{2}\right) \cdot \zeta\left(\frac{s}{2} - 3\right)$$
		Die meromorph fortgesetzte Funktion $Z_{\mathbf{E}_8}(s)$ besitzt in $\mathbb{C}$ exakt eine Polstelle erster Ordnung bei $s = 8$. Das Residuum beträgt:
		$$\operatorname*{Res}_{s=8} Z_{\mathbf{E}_8}(s) = \frac{\pi^4}{3}$$
		Die scheinbare Singularität bei $s = 2$ ist hebbar; der Wert der analytischen Fortsetzung lautet dort:
		$$Z_{\mathbf{E}_8}(2) = -\frac{30\,\zeta(3)}{\pi^2}$$
	\end{theorem}
	
	\begin{proof}
		Die Faktorisierung folgt aus $|S_n| = 240 \sigma_3(n)$ und der Dirichlet-Identität $\sum_{n=1}^\infty \sigma_3(n) n^{-u} = \zeta(u)\zeta(u-3)$ mit $u = s/2$ \cite{apostol1976}.
		Für $u \to 4$ ($s \to 8$) gilt $\zeta(u-3) = \frac{1}{u-4} + \mathcal{O}(1) = \frac{2}{s-8} + \mathcal{O}(1)$. Das Residuum berechnet sich zu $240 \cdot 2^{-4} \cdot \zeta(4) \cdot 2 = 30 \cdot \frac{\pi^4}{90} = \frac{\pi^4}{3}$.
		Für $u \to 1$ ($s \to 2$) gilt $\zeta(u) = \frac{1}{u-1} + \mathcal{O}(1)$ und $\zeta(u-3) = \zeta(-2 + (u-1)) = \zeta'(-2)(u-1) + \mathcal{O}((u-1)^2)$. Pol und Nullstelle heben sich auf. Mit der bekannten Identität $\zeta'(-2) = -\frac{\zeta(3)}{4\pi^2}$ folgt der angegebene Fortsetzungswert $Z_{\mathbf{E}_8}(2) = 120\,\zeta'(-2) = -\frac{30\zeta(3)}{\pi^2}$.
	\end{proof}
	
	\begin{remark}
		Da die ursprüngliche Dirichlet-Reihe nur für $\operatorname{Re}(s) > 8$ konvergiert, ist der negative Wert $Z_{\mathbf{E}_8}(2) < 0$ als Funktionswert der analytischen Fortsetzung und nicht als Summe positiver Reihenterme zu verstehen.
	\end{remark}
	
	\begin{theorem}[Skalierungszerlegung und Möbius-Inversion für primitive Vektoren]
		Sei $\mathbf{E}_8^{\text{prim}} := \{x \in \mathbf{E}_8 \setminus \{0\} \mid x \notin m\mathbf{E}_8 \text{ für alle } m \ge 2\}$ die Menge der primitiven Gittervektoren. Für $\operatorname{Re}(s) > 8$ gilt für die primitive Epstein-Zetafunktion:
		$$Z_{\mathbf{E}_8}^{\text{prim}}(s) := \sum_{x \in \mathbf{E}_8^{\text{prim}}} \|x\|^{-s} = \frac{Z_{\mathbf{E}_8}(s)}{\zeta(s)} = 240 \cdot 2^{-s/2} \cdot \frac{\zeta(s/2)\,\zeta(s/2 - 3)}{\zeta(s)}$$
		Die rechte Seite liefert eine meromorphe Fortsetzung der primitiven Zetafunktion.
	\end{theorem}
	
	\begin{proof}
		Aus der eindeutigen Gitterzerlegung $x = m \cdot y$ mit $m \in \mathbb{N}$ und $y \in \mathbf{E}_8^{\text{prim}}$ folgt wegen $\|x\|^{-s} = m^{-s} \|y\|^{-s}$ im Konvergenzbereich $\operatorname{Re}(s) > 8$ direkt $Z_{\mathbf{E}_8}(s) = \sum_{m=1}^\infty m^{-s} Z_{\mathbf{E}_8}^{\text{prim}}(s) = \zeta(s) Z_{\mathbf{E}_8}^{\text{prim}}(s)$. Die Division durch $\zeta(s)$ entspricht der Dirichletschen Möbius-Inversion wegen $\zeta(s)^{-1} = \sum_{m=1}^\infty \mu(m) m^{-s}$.
	\end{proof}
	
	\begin{remark}
		Anders als $Z_{\mathbf{E}_8}(s)$ kann die meromorphe Fortsetzung der primitiven Funktion $Z_{\mathbf{E}_8}^{\text{prim}}(s)$ aufgrund des Faktors $1/\zeta(s)$ zusätzliche Pole an den nicht-trivialen Nullstellen der Riemannschen Zeta-Funktion besitzen, soweit keine Aufhebung mit dem Zähler eintritt.
	\end{remark}
	
	\section{Fazit}
	Die radiale Initialtopologie $(\mathbf{E}_8, \tau_{\text{rad}})$ organisiert die Normschalen als topologisch ununterscheidbare Fasern der Normabbildung $Q$. Der zugehörige Kolmogorov-Quotient identifiziert die Schalenstruktur kanonisch mit der diskreten Menge $\mathbb{N}_0$. Die gewichtete Schalen-Zählfunktion führt auf die klassische $E_8$-Theta-Reihe und deren meromorphe Dirichlet-Fortsetzung mit einem einzigen Pol bei $s=8$.
	
	\begin{thebibliography}{9}
		
		\bibitem{serre1973}
		J.-P. Serre, \emph{A Course in Arithmetic}, Graduate Texts in Mathematics, vol.~7, Springer-Verlag, New York, 1973.
		
		\bibitem{conway1999}
		J. H. Conway and N. J. A. Sloane, \emph{Sphere Packings, Lattices and Groups}, 3rd ed., Grundlehren der mathematischen Wissenschaften, vol.~290, Springer-Verlag, New York, 1999.
		
		\bibitem{apostol1976}
		T. M. Apostol, \emph{Introduction to Analytic Number Theory}, Undergraduate Texts in Mathematics, Springer-Verlag, New York, 1976.
		
		\bibitem{terras1985}
		A. Terras, \emph{Harmonic Analysis on Symmetric Spaces and Applications I}, Springer-Verlag, New York, 1985.
		
		\bibitem{viazovska2017}
		M. S. Viazovska, \emph{The sphere packing problem in dimension 8}, Annals of Mathematics \textbf{185} (2017), no.~3, 991--1015.
		
	\end{thebibliography}
	
\end{document}